% Generated draft: requires mathematical and source review.
\documentclass[11pt]{article}
\usepackage[margin=25mm]{geometry}
\usepackage{fontspec}
\setmainfont{Latin Modern Roman}
\usepackage{amsmath,amssymb,mathtools,unicode-math}
\setmathfont{Latin Modern Math}
\usepackage{enumitem,needspace,xurl,hyperref}
\setlength{\parindent}{0pt}
\setlength{\parskip}{5pt}
\setlength{\emergencystretch}{3em}
\newcommand{\N}{\mathbb{N}}
\newcommand{\Z}{\mathbb{Z}}
\newcommand{\Q}{\mathbb{Q}}
\newcommand{\R}{\mathbb{R}}
\newcommand{\C}{\mathbb{C}}
\newcommand{\F}{\mathbb{F}}
\newcommand{\A}{\mathbb{A}}
\newcommand{\PP}{\mathbb{P}}
\newcommand{\E}{\mathbb{E}}
\newcommand{\eps}{\varepsilon}
\newcommand{\dd}{\,\mathrm{d}}
\begin{document}
\section*{UnsolvedMath ID: 8700032}
\textbf{Source problem code:} catalogue entry
\subsection{Definitions and mathematical statement}
Let $n\ge1$, let $\Sigma\subset\C^n$ be finite, and let $\varepsilon>0$. An entire function $f:\C^n\to\C$ vanishes on $\Sigma$ with multiplicity at least $t$ if
\[
 \partial^\alpha f(z)=0\qquad
 (z\in\Sigma,\ \alpha\in\N_0^n,\ |\alpha|<t),
\]
where $|\alpha|=\sum_i\alpha_i$.

The conjecture asserts the uniform quantifier order
\[
 \exists r_0=r_0(\Sigma,\varepsilon)>0\quad
 \forall t\in\Z_{>0}\quad
 \forall f\text{ entire with the stated vanishing}\quad
 \forall r\ge r_0,
 \qquad
 \Theta_f(r)\ge\omega_t(\Sigma)-t\varepsilon.
\]
In particular the radius threshold is allowed to depend on $\Sigma$ and $\varepsilon$, but not on $t$ or $f$. There is no stated restriction to polynomial functions, finite-order entire functions, or a particular growth class.

The source does not define $\Theta_f(r)$ or $\omega_t(\Sigma)$. The summary only describes them as growth/counting quantities and finite-set exponents, which does not determine their normalization or formula. They are retained as source symbols rather than replaced by a guessed Nevanlinna characteristic or polynomial-degree invariant. The treatment of the identically zero function, the empty finite set, and any conventions required to make those quantities defined must also be supplied; the source does not exclude those cases.
\subsection{Short English statement}
A single radius threshold should make the stated growth lower bound hold for every multiplicity and every entire function vanishing to that multiplicity on the fixed finite set. Uniformity in both the function and multiplicity is central.

The supplied summary describes the assertion as a tool-oriented survey conjecture related to value distribution and transcendence measures, with weaker surrounding estimates, but gives no precise estimate or isolated published resolution. The OPEN-TRIAGE label and 2004 survey are retained with current primary-literature status audit still required.
\subsection{Sources}
\begin{itemize}
\item Local UnsolvedMath registry, permanent ID 8700032, source code catalogue entry, ``Conjecture 3.14,'' proposer metadata Michel Waldschmidt. Dataset curation and metadata: Ulam AI and UnsolvedMath Contributors, CC BY 4.0.
\item Michel Waldschmidt, \emph{Open Diophantine Problems} (2004), Conjecture 3.14, \url{https://arxiv.org/abs/math/0312440}; source-TeX extraction, accessed 2026-07-29.
\item No separate primary resolution or precise weaker estimate is supplied. Source-specific provenance: the survey presents the item as open; current status audit and redistribution terms NEEDS\_REVIEW for the publicly accessible author/publisher copy, separately from the dataset licence.
\end{itemize}
\end{document}
